\documentclass[a4paper, oneside]{article}
\usepackage{graphicx}
\usepackage{amsthm}
\usepackage{amsmath}
\usepackage{amssymb}
\usepackage[a4paper,
            bindingoffset=0.2in,
            left=2cm,
            right=2cm,
            top=2cm,
            bottom=2cm,
            footskip=.25in]{geometry}
\usepackage[italian]{babel}
\usepackage{pgfplots}
\usepackage{tabularx}
\usepackage{tikz}
\usepackage{wrapfig}
\usepackage{color}
\usepackage[d]{esvect}
\definecolor{page}{rgb}{0.129,0.157,0.212}
\pagecolor{page}
\color{white}
\graphicspath{ {./images/} }
\usetikzlibrary{shapes.geometric}
\usetikzlibrary{datavisualization}
\usetikzlibrary{datavisualization.formats.functions}
\usetikzlibrary{patterns}
\pgfplotsset{width=10cm,compat=1.18}

\title{Appunti di MEccanica}
\author{Tommaso Miliani}
\date{10-03-26}

\begin{document}
\newtheoremstyle{theoremEnv}
                {}          % Space above
                {}          % Space below
                {\slshape}  % Body font
                {}          % Indent amount
                {\bfseries} % Head font
                {.}         % Punctuation after head
                {\newline}  % Space after theorem head
                {}          % Theorem head spec
\theoremstyle{theoremEnv}

\newtheorem{definition}{Definizione}[section]
\newtheorem{theorem}{Teorema}[section]
\newtheorem{lemma}{Proposizione}[section]
\newtheorem{observation}{Osservazione}[section]
\newtheorem{corollary}{Corollario}[theorem]
\newtheorem{example}{Esempio}[section]
\newtheorem{remark}{Enunciato}[section]

\maketitle

\section{Curve parametriche}
\begin{wrapfigure}{r}{0.4\textwidth}
    \centering
    \caption{}
    \begin{tikzpicture}
        \draw(0, 0) arc (120:60:1);
    \end{tikzpicture}    
\end{wrapfigure}
Data una parametrizzazione $\vv{r}(t)$, si definisce il \textbf{vettore curvatura} 
come
\begin{gather*}
    \vv{k} = \frac{d\hat{t}  }{ds}  \qquad \hat{t} = \frac{d\vv{r} }{ds} 
\end{gather*}
Si chiama \textbf{curvatura} il modulo di tale vettore. Si osserva come 
il vettore tangente ed il vettore curvatura siano due vettori ortogonali tra di loro (infatti 
la derivata di un versore è ortogonale al versore di partenza).
\begin{example}[Calcolo della curvatura di un cerchio]
    Preso un cerchio di raggio $R$
    \begin{gather*}
        \begin{tikzpicture}
            \draw(-1, 0) arc (180:0:1);
            \draw[->](-1.5, 0) -- (1.5, 0); 
            \draw[->](0, 0) -- (0, 1.5) node[at end, left] {$\hat{y}$ };
            \draw(0, 0) -- (0.84, 0.5) node[midway, above] {$R$};
            \filldraw(0, 0) node[anchor = north] {$O$};
            \filldraw(0.84, 0.5) node[anchor = west] {$A$};
        \end{tikzpicture}
    \end{gather*}
    SI può esprimere, con la parametrizzazione circolare, gli angoli sul cerchio come 
    \begin{gather*}
        \vv{r}(\phi)  = (A - O) = R(\cos\phi \hat{x} + \sin\phi \hat{y}  )
    \end{gather*}
    Si può allora ottenere 
    \begin{gather*}
        s(\phi) = \int_{0}^{\phi} \left| \frac{d\vv{r} }{d\phi}  \right| \ \phi  \qquad \frac{d\vv{r} }{ds} =  R(-\sin\phi \hat{x} + \hat{j} \cos\phi  ) 
    \end{gather*}
    Allora l'integrale diventa
    \begin{gather*}
        s(\phi) = \int_{0}^{\phi} R \ d\phi = R\int_{0}^{\phi} d\phi = R\phi  
    \end{gather*}
    Per esprimere $R$ in funzione di $s$: 
    \begin{gather*}
        \phi = \frac{S}{R} \ \Longrightarrow \ \vv{r}(S) = R(\hat{x} \cos\frac{S}{R} + \hat{y} \sin\frac{S}{R}  ) 
    \end{gather*}
    Per ottenere la curvatura, si esegue la derivata seconda di $\vv{r}$ (ossia
    di quello che si è appena trovato):
    \begin{align*}
        \frac{d\vv{r} }{dS} &= R\left(-\hat{i}\sin\left(\frac{S}{R}\right) \frac{1}{R} + \hat{y} \cos\left(\frac{S}{R}\right) \frac{1}{R}  \right) \\
        \frac{d\vv{r} }{dS^{2}} &= -\hat{x}\cos\left(\frac{S}{R}\right) \frac{1}{R} + -\hat{y}\sin\left(\frac{S}{R} \right) \frac{1}{R}    = - \frac{\vv{r}(S) }{R^{2}}
    \end{align*} 
    Allora si ottiene la curvatura come 
    \begin{gather*}
        k = \left| \frac{d^{2}\vv{r} }{dS^{2}} \right|  = \frac{1}{R}
    \end{gather*}
\end{example}

\begin{theorem}
    Sia $r(t)$ una curva in $C^{2}(\mathbb{R})$ una curva regolare con parametro $t \in \mathbb{R}$. Si ha che
    \begin{gather*}
        \frac{d^{2}r}{dS^{2}} = \left(\frac{d^{2}r}{dt^{2}} - \frac{\left< \frac{d^{2}r}{dS^{2}}, \frac{d^{2}r}{dt^{2}} \right> }{\left| \frac{d^{2}r}{dt^{2}} \right|^{2} } \frac{dr}{dt}\right) \frac{1}{\left| \frac{dr}{dt} \right|^{2} }
    \end{gather*}
    E la \textbf{curvatura}:
    \begin{gather*}
        k = \frac{\left| \frac{d\vv{r} }{dt} \right| \wedge \frac{d^{2}r}{dt^{2}} }{\left| \frac{dr}{dt} \right|^{3} }
    \end{gather*}
\end{theorem}
\begin{proof}
    Si esprimere la derivata della funzione composta come
    \begin{gather*}
        \frac{d\vv{r} }{dS} = \frac{d\vv{r} }{dt} \frac{dt}{dS} = \frac{1}{\left|  \frac{d\vv{r} }{dt}\right| } \frac{d\vv{r} }{dt} \\
        \frac{d^{2}r}{dS^{2}} = \frac{d}{dS} \frac{dr}{dS} = \frac{d}{dS} \left(\dot{r} \left| \dot{r}  \right|^{-1}  \right) = \frac{d}{dt} \left(\dot{r} \left| \dot{r}  \right|^{-1}  \right) \frac{dt}{dS} = \frac{1}{\left| \dot{r}  \right| } \frac{d}{dt} (\dot{r} \left| \dot{r}  \right| ^{-1}  )
    \end{gather*}
    DUnque
    \begin{gather*}
        \frac{1}{\left| \dot{r}  \right| } \left(\ddot{r} \left| \dot{r}  \right| ^{-1} + \dot{r} \frac{d}{dt}(\left| \dot{r}  \right| ^{-1} )   \right)
    \end{gather*}
    L'ultima derivata diventa:
    \begin{gather*}
        \frac{d}{dt}(\left| \dot{r}  \right| ^{-1} ) = \frac{d}{dt} (\left< \dot{r}, \dot{r}   \right> ^{-1/2} ) = - \frac{1}{2} \left< \dot{r}, \dot{r}   \right>^{-3/2} \frac{d}{dt}\left< \dot{r}, \dot{r}   \right> = -\frac{1}{2}  \left< \dot{r}, \dot{r}   \right>^{-3/2} (\left< \ddot{r}, \dot{r}   \right> + \left< \dot{r}, \ddot{r }   \right>  ) = -\left< \dot{r}, \ddot{r}   \right> \left| \dot{r}  \right|^{-3} 
    \end{gather*}
    Allora 
    \begin{gather*}
        \frac{d^{2}r}{dS^{2}}\frac{1}{\left| \dot{r}  \right| } \left(\ddot{r} \left| \dot{r}  \right| ^{-1}  - \dot{r} \frac{\left< \dot{r}, \ddot{r}   \right> }{\left| \dot{r}  \right|^3 }  \right) = \frac{1}{\left| \dot{r}  \right|^2 } \left(\ddot{r} + \dot{r} \frac{\left< \dot{r}, \ddot{r}   \right> }{\left| \dot{r}  \right|^2 }  \right)
    \end{gather*}
    Adesso si esegue il prodotto vettoriale a con l'espressione appena trovata
    rispetto al vettore $\frac{d\vv{r} }{dS}$:
    \begin{gather*}
        \frac{d^{2}r}{ds^{2}} \wedge \frac{d\vv{r} }{ds} = \frac{1}{\left| \dot{r}  \right|^{3} }\left(\ddot{r} \wedge \dot{r}  - \frac{\left< \dot{r}, \ddot{r}   \right>}{\left| \dot{r}  \right| } \dot{r} \wedge \dot{r}   \right)
    \end{gather*}
    Si deve dimostrare che questa espressione sia esattamente la curvatura. Dunque si calcola il modulo 
    dell'espressione ottenuta ora:
    \begin{gather*}
        \sin\left(\frac{d^{2}r}{ds^{2}} \angle \frac{dr}{ds}\right) \left| \frac{d^{2}r}{ds^{2}} \right| \left| \frac{dr}{ds} \right|  
    \end{gather*}
    Il primo modulo corrisponde al modulo del vettore curvatura ed il secondo è uno, dunque l'angolo tra i due è
    90 e allora il modulo è esattamente il modulo del vettore curvatura.
\end{proof}

\begin{example}
    Data una funzione regolare $f(x)  : \mathbb{R} \to \mathbb{R}$ che sia derivabile 
    in tutto il suo dominio (è una curva piana).
    \begin{gather*}
        y = f(x) \qquad \vv{r}(x) = x\hat{x} + f(x) \hat{y}   
    \end{gather*}
    Si vuole determinare la curvatura di questa funzione:
    \begin{gather*}
        \frac{d\vv{r} }{dx} = \hat{x} + f'(x) \hat{y} \qquad \left| \frac{dr}{dx} \right| = \sqrt{1 + (f'(x))^{2}}    
    \end{gather*}
    La curva è regolare se è $C^{1}$ poiché l'argomento della radice è sempre positivo. La derivata seconda
    \begin{gather*}
        \frac{d^{2}\vv{r} }{dx^{2}} = f''(x)\hat{y} 
    \end{gather*} 
    Dunque la curvatura:
    \begin{gather*}
        k = \frac{1}{(1 + (f'(x))^{2})^{1/2}} \left| (\hat{x} + \hat{y}f(x)   ) \wedge   f''(x) \hat{y} \right|  = \frac{\left| f'(x) \right| }{(1 + (f'(x))^{2})^{1/2}}
    \end{gather*}
\end{example}

\begin{example}
    Determinare il vettoe
    re tangente e la curva tangente al piano $r(1)$:
    \begin{gather*}
        \vv{r}(t) = (1 + t)\hat{x} - t^{2} \hat{j} + (1 + t^{3}) \hat{z} \qquad \vv{r}(1) = 2\hat{x} - \hat{y} + 2\hat{z}       
    \end{gather*}
    La retta tangente sarà data da:
    \begin{gather*}
        \vv{v}  = \vv{r}(1) + \alpha \frac{d\vv{r} }{dt} 
    \end{gather*}
    La quale va calcolata in $t = 1$:
    \begin{gather*}
        \frac{d\vv{r} }{dt} = \hat{x} - 2t\hat{y}   + 2t^{2}\hat{z} = \hat{x} - 2\hat{y} + 2\hat{z}   
    \end{gather*}
    Questo mi permette di ottenere la retta:
    \begin{gather*}
        \vv{v} =  (2\hat{x} - \hat{y} + 2\hat{z }   ) + \alpha (\hat{x} - 2\hat{y} + 2\hat{z}   )
    \end{gather*}
    Il vettore tangente è anche il vettore perpendicolare al piano considerato 
\end{example}

\section{Superfici in $\mathbb{R}^{3}$}
Si studiano le proprietà fondamentali delle superfici in $\mathbb{R}^{3}$ che ci consentono di 
studiare le \textbf{ipersuperfici}. Una superficie è una deformazione continua e 
regolare di un pezzo di piano. Ci sono vari modi per poter definire le superfici, ma noi la si vede
come una mappa lineare. Preso un $\mathbb{U} \subset \mathbb{R}^{2}$:
\begin{align}
    \phi : \mathbb{R}^{2} \to \mathbb{R}^{3}
\end{align}
Quindi si definisce la \textbf{superficie} $\Sigma \equiv I_m(\phi(\mathbb{U}))$. Se si volesse 
scrivere $\phi$ in maniera esplicita, si potrebbe scrivere
\begin{align}
    \vv{\phi}(u_1, u_2) = \begin{pmatrix} x(u_1, u_2)  \\
    y(u_1, u_2) \\
    z(u_1, u_2) \end{pmatrix}   = \hat{x} x(u_1, u_2) + \hat{y} y(u_1, u_2) + \hat{z} z(u_1, u_2)   
\end{align} 
Si chiede inoltre che $\vv{\phi} \in C^{1}(\mathbb{U})$, in modo tale che le i-esime 
componenti siano delle funzioni continue e derivabili. A differenza delle curve le superfici è molto difficile (il 
passaggio da uno a due parametri è il più difficile). Per poter studiare la superficie si studiano le curve 
che compongono la superficie: studiando le proprietà delle curve si possono ricavare 
le proprietà della curva stessa. Preso un punto $A$ nello spazio dei parametri, 
esso corrisponderà ad un certo punto $B$ secondo la mappa $\phi$. L'idea è dunque costrutire delle
curve semplici che passino per $B$ (la curva più semplice è l'immagine di un vettore parallelo all'asse $x$).  \\
Si costruisce una curva orizzontale che passa per il punto $A$ che prende il nome 
di $\gamma(t)$ la quale è, formalmente, 
\begin{gather*}
    \gamma(t) = (A - O) + \hat{x}t = \begin{pmatrix} u_1 + t \\
    u_2 \end{pmatrix} 
\end{gather*}
La cui curva corrispondente sulla superficie è
\begin{gather*}
    r(\gamma(t)) = \vv{r} (u_1 + t, u_2)  
\end{gather*}
Dunque, dipendendo da un solo parametro, si può ottenere la tangente a questa curva:
\begin{gather*}
    \frac{d\vv{r} }{dt} = \lim_{t \to 0} \frac{\vv{r}(u_1 + t, u_2) - \vv{r}(u_1 , u_2)  }{t} = \frac{\partial \vv{r} }{\partial u_1}   = \hat{x} \frac{\partial x}{\partial u_1} + \hat{y} \frac{\partial y}{\partial u_1} + \hat{z}\frac{\partial z}{\partial u_1}      
\end{gather*}
Ossia la derivata direzionale della funzione rispetto a $u_1$.  Si ottengono
allora i \textbf{vettori coordinati} (o \textbf{vettori coordinate}):
\begin{gather*}
    \frac{\partial \vv{r} }{\partial u_1} \quad \frac{\partial \vv{r} }{\partial u_2}  
\end{gather*}
Che in generale non sono né ortogonali né paralleli tra loro ma possono essere vettori generici 
di cui non si conosce bene la relazione tra loro. 
Si vuole ora studiare il piano generato dallo span
di $\vv{a}, \vv{b} \in \mathbb{R}^{3}$, con $\alpha, \beta \in \mathbb{R}$ si generano i seguenti vettori:
\begin{gather*}
    v = \alpha\vv{a} + \beta\vv{b} 
\end{gather*}  
Dunque una superficie può essere scritta come 
\begin{gather*}
    \vv{r}= u_1 \vv{a} + u_2 \vv{b}   + (x_0 \hat{x} + y_0 \hat{j} + z_0 \hat{z}   )
\end{gather*}
Ossia il piano che si ottiene traslato rispetto all'origine:
\begin{gather*}
    (P_0 - O) = x_0 \hat{x} + y_0 \hat{y} + z_0 \hat{z}   
\end{gather*}
Allora 
\begin{gather*}
    \frac{\partial \vv{r} }{\partial u_1 } = \vv{a} \qquad \frac{\partial \vv{r} }{\partial u_2} = \vv{b}    
\end{gather*}

\begin{example}[Vettori coordinati di una superficie seferica]
    Presa una sfera di raggio $R$:
    \begin{gather*}
        \begin{tikzpicture}
            \draw(0, 0) circle (1);
        \end{tikzpicture}
    \end{gather*}
    Utilizzando le coordinate sferiche, si prende un generico vettore
    parametrizzato secondo le coordinate sferiche, in modo tale che 
    le parametrizzazioni siano $\phi \in [0, 2\pi]$, $\theta \in [0, \pi]$.
    Preso un punto qualsiasi si hanno due possibilità: ci si può muovere o a 
    $\theta$ costante o a $\phi$ costante. 
    \begin{gather*}
        x = R\sin\theta \cos\phi \\
        y = R \sin\theta \sin\phi  \\
        z = R\sin\theta 
    \end{gather*}
    Derivando rispetto a $\theta$:
    \begin{gather*}
        \frac{\partial \vv{r} }{\partial \theta} = R\left(\hat{x}\cos\theta \cos\phi + \hat{y}\cos\theta\sin\phi - \hat{z}\sin\theta   \right) 
    \end{gather*}
    Mentre la derivata rispetto a $\phi$:
    \begin{gather*}
        \frac{\partial \vv{r} }{\partial \phi} = R\left(-\sin\theta \sin\phi \hat{x} + \sin\theta \cos\phi \hat{y}  \right) 
    \end{gather*}
    I due vettori risultano ortogonali tra loro (come volevasi vedere). 
\end{example}

\section{Piano tangente (ad una superficie)}
Il piano tangente è il prototipo di uno \textbf{spazio tangente}, 
ossia spazi vettoriali tangenti ad altri spazi vettoriali. Il piano tangente è, 
come lo spazio vettoriale che continene tutti i vettori tangenti ad una data superficie passanti
per un determinato punto. Se la curva è arbitraria allora sto considerando tutti i vettori 
che hanno questa caratteristica. 
\begin{align}
    \mathcal{T}_{r_0} (\Sigma) = \{ \vv{v} \in \mathbb{R}^{3} : \exists t_0 \in \mathbb{R} \ : \ \vv{\phi}(t) : \mathbb{R} \to \Sigma  \ t.c. \ v = \frac{d\vv{r}(t_0) }{dt} \ e \ \vv{r} (t_0) = r_0\}
\end{align}
Che non è scontato che sia spazio vettoriale.\\
Sia $\Sigma = Im(\phi(\mathbb{U}))$ tale che $\phi :\mathbb{U} \subset \mathbb{R}^{2} \to \mathbb{R}^{3}$ e $\phi \in C^{1}(\mathbb{U})$. $\phi$ è 
la \textbf{parametrizzazione} della superficie e si definisce \textbf{regolare} se $\phi \in C^{1}(\mathbb{U})$ e 
\begin{gather*}
    \frac{\partial \vv{\phi} }{\partial u_1} \qquad \frac{\partial \vv{\phi} }{\partial u_2}  
\end{gather*}
Sono linearmente indipendenti. Allora questa prende il nome di \textbf{parametrizzazione regolare}. Si dirà, 
inoltre, che $\Sigma$ è regolare se ammette parametrizzazione regolare. 

\begin{theorem}
    I vettori coordinati della parametrizzazione sono linearmente indipendenti se 
    la Jacobiana di $\phi$ ha rango massimo (2). Si ricorda che la matrice Jacobiana è
    la matrice che 
    \begin{gather*}
        D_{u_1, u_2}\phi  = \begin{pmatrix}
            \dfrac{\partial \phi_x}{\partial u_1} & \dfrac{\partial \phi_x}{\partial u_2} \\
             \dfrac{\partial \phi_y}{\partial u_1} & \dfrac{\partial \phi_y}{\partial u_2} \\
            \dfrac{\partial \phi_z}{\partial u_1} & \dfrac{\partial \phi_z}{\partial u_2} \\
        \end{pmatrix} = \begin{pmatrix}
            \frac{\partial \vv{\phi} }{\partial u_1} & \frac{\partial \vv{\phi} }{\partial u_2}  
        \end{pmatrix}
    \end{gather*}
    Ossia la matriche che ha come colonne le derivate delle componenti di $\phi$ rispetto ai 
    parametri.
\end{theorem}

\end{document}